% TOP_PROBLEM: {"id": 30006635, "title": "Liouville Theory from the Virasoro Conformal Net", "category_id": 16, "category": {"id": 16, "name": "physics", "display_name": "Mathematical Physics", "description": "Problems at the intersection of mathematics and physics.", "slug": "physics", "order_index": 16, "created_at": "2026-07-31T15:26:25.671Z"}, "status": "open", "statusQualification": "For a nonrational net the source construction is defined on the subcategory whose connected surface components have nonempty incoming and outgoing boundaries. An extension to every conformal bordism is not silently assumed."}
% FIRST_POSTED: 2026-10-11
% ENTRY_KIND: statement_only
% SUBMISSION: {"kind": "statement_only", "category": "top_problems", "problem_id": "30006635", "proposal_type": "established", "source": "Higher Structures from Symmetries in Quantum Field Theory (2026), Oberwolfach Reports, DOI 10.4171/owr/2026/12, pp. 737--794. Original report and mathematical context. Incorporated report material remains under the source CC BY-SA 4.0 terms unless otherwise noted. https://doi.org/10.4171/owr/2026/12", "status_evidence": "Source status retained from the supplied canonical registry. Dated literature assessment (2026-08-22): The Virasoro-net construction is still presented as conjecturally producing Liouville theory; recent work proves a net-to-unitary-VOA direction but not the claimed identification with Liouville theory. Primary report inspected for formulation; no comprehensive current-status literature review is claimed.", "status_checked_on": "2026-10-11", "duplicate_check": "Canonical numeric ID and public code checked against full registry; absent from baseline selected website registry. Source scope compared with adjacent source questions and existing related catalogue entries. No new identity allocated.", "attribution": "Original authors and contributors credited in Higher Structures from Symmetries in Quantum Field Theory (2026), Oberwolfach Reports, DOI 10.4171/owr/2026/12, pp. 737--794; UnsolvedMath curation credited separately. Definitions and English exposition prepared with Codex assistance; no independent human mathematical verification claimed."}
% SOURCE_STATEMENT: "Prove that applying the fully extended two-dimensional conformal-field-theory construction to the Virasoro conformal net produces Liouville conformal field theory."
% Preparation and final self-audit: CONTRIBUTING.md prompts C and D; source/audit artifacts retained in the local physics_statements_2026_10_11 batch.
% Imported from: physics_statements.tex; UnsolvedMath ID 30006635
% !TeX program = lualatex
\documentclass[11pt]{article}
\usepackage[margin=25mm]{geometry}
\usepackage{fontspec}
\setmainfont{Latin Modern Roman}
\usepackage{amsmath,amssymb,mathtools}
\usepackage{unicode-math}
\setmathfont{Latin Modern Math}
\usepackage{xurl,hyperref}
\hypersetup{colorlinks=true,urlcolor=blue}
\setcounter{secnumdepth}{0}
\setlength{\parindent}{0pt}
\setlength{\parskip}{5pt}
\setlength{\emergencystretch}{3em}
\newcommand{\sref}[2]{\hyperlink{source:#1:#2}{[S#2]}}
\newcommand{\cataloguescope}[1]{\paragraph{Recorded target.}#1}
\newcommand{\R}{\mathbb R}
\newcommand{\C}{\mathbb C}
\newcommand{\N}{\mathbb N}
\newcommand{\Z}{\mathbb Z}
\newcommand{\Q}{\mathbb Q}
\newcommand{\F}{\mathbb F}
\newcommand{\eps}{\varepsilon}
\begin{document}
% UnsolvedMath ID 30006635; source catalogue code OWR-14299913-004
\setcounter{section}{30006634}
\section{Liouville Theory from the Virasoro Conformal Net}\label{30006635}
{\small UnsolvedMath ID 30006635 · Mathematical Physics}
\subsection{Definitions and mathematical statement}

\paragraph{Shared notation.}$\mathbb N=\{1,2,\ldots\}$; the standard number fields have their usual meanings. All additional conventions are stated locally.


A conformal net assigns local von Neumann algebras to intervals, with locality and conformal covariance. The Virasoro net is generated by the corresponding stress-energy representation. The reported construction assigns algebras, bimodules and intertwiners to points, one-dimensional bordisms and conformal surfaces. Liouville theory is the proposed resulting nonrational conformal theory.

\paragraph{Catalogue wording retained for traceability.}
Prove that applying the fully extended two-dimensional conformal-field-theory construction to the Virasoro conformal net produces Liouville conformal field theory.

\paragraph{Scope and source qualification.}For a nonrational net the source construction is defined on the subcategory whose connected surface components have nonempty incoming and outgoing boundaries. An extension to every conformal bordism is not silently assumed.
\subsection{Short English statement}
Identify the field theory produced by the source conformal-net construction with Liouville theory.
\subsection{Sources}
\begin{itemize}
\item[\hypertarget{source:30006635:1}{S1}] ULAM.AI and the UnsolvedMath Contributors, UnsolvedMath, record 30006635 (OWR-14299913-004). Catalogue curation under CC BY 4.0; definitions and English exposition added locally.
\url{https://huggingface.co/datasets/ulamai/UnsolvedMath}.
\item[\hypertarget{source:30006635:2}{S2}] Higher Structures from Symmetries in Quantum Field Theory (2026), Oberwolfach Reports, DOI 10.4171/owr/2026/12, pp. 737--794. Original report and mathematical context. Incorporated report material remains under the source CC BY-SA 4.0 terms unless otherwise noted.
\url{https://doi.org/10.4171/owr/2026/12}.
\end{itemize}
\label{end:30006635}
\end{document}
